\section{Reproduction}

\textbf{Artifact.} The code, the request traces and the result files
behind every figure and table are packaged as
\texttt{talg-artifact.zip}, released at

\begin{quote}
\texttt{https://github.com/bruce5800/Postgraduate-Project/releases/tag/talg-submission-2026-08}
\end{quote}

Unpack it and run the commands below from its root; \texttt{README.md}
inside repeats the essentials. The artifact also contains
\texttt{budgetstakes/}, the Lean 4 development of Appendix C
(\texttt{lake\ build} replays the check). \texttt{results/} holds our
own outputs, so a fresh run can be diffed against them directly.

Every number in this paper is regenerated from a fixed seed by a single
script. Almost everything runs on a clean checkout of the artifact: the
synthetic experiments generate their own instances, and the request
traces behind Sections 6.1, 6.3 and 8 are small enough that we ship them
with the artifact under \texttt{data/trace/}. One class is the exception
--- the six real graphs of Section 6.2 (and the Phase-2 reproduction
that reuses them) are third-party datasets we do not redistribute, so
they must be downloaded into \texttt{data/realworld/} first. A dagger
(\(\dagger\)) marks that class, and only that class, in the map below;
§A.4 says what each dataset is and where it comes from.

\subsection{Environment and
conventions}

\begin{itemize}
\tightlist
\item
  \textbf{Stack:} Python 3.12; NumPy 1.26, SciPy 1.13, NetworkX 3.3,
  Matplotlib (plots only).
\item
  \textbf{Randomness:} every stream derives from one master seed
  (default \texttt{0}), spawned into independent reproducible
  sub-streams with NumPy's \texttt{default\_rng(seed).spawn(k)} --- one
  each for the graph, the arrival sequence, algorithm randomness and the
  perturbation.
\item
  \textbf{Paired trials:} within a comparison every algorithm and every
  error level reuses the same graph, arrival sequence, \(\mathrm{OPT}\)
  and tie-break seed.
\item
  \textbf{Confidence intervals:} 95\% normal-approximation half-widths
  over trials.
\item
  \textbf{Flow decompositions.} Feldman, Jaillet--Lu and the serving
  advice \(b\)-matching consume the \emph{decomposition} of a maximum
  flow, which --- unlike the flow value --- is not unique. Their
  networks label nodes with integers rather than strings, so the
  decomposition NetworkX returns is stable across runs; under string
  labels it followed the interpreter's per-process string hashing and
  every derived number moved between runs of the same script (Section
  2.5).
\item
  Runtimes below are wall-clock on one machine (Apple M4 Pro, 12 cores).
  The scripts are single-threaded apart from NumPy's own vectorization,
  so they scale with single-core speed.
\end{itemize}

\subsection{\texorpdfstring{Result \ensuremath{\rightarrow} script
map}{Result  script map}}

Each row names a \emph{stem}. Unless the row says otherwise, its script
is \texttt{run\_\textless{}stem\textgreater{}.py} under
\texttt{scripts/}, and its outputs are
\texttt{\textless{}stem\textgreater{}.json} and
\texttt{\textless{}stem\textgreater{}.png} under \texttt{results/}.

{ % do not increment counter
\begin{longtable}[]{@{}
  >{\raggedright\arraybackslash}p{(\linewidth - 4\tabcolsep) * \real{0.3333}}
  >{\raggedright\arraybackslash}p{(\linewidth - 4\tabcolsep) * \real{0.3333}}
  >{\raggedright\arraybackslash}p{(\linewidth - 4\tabcolsep) * \real{0.3333}}@{}}
\toprule\noalign{}
\begin{minipage}[b]{\linewidth}\raggedright
Paper object
\end{minipage} & \begin{minipage}[b]{\linewidth}\raggedright
Script stem (or path)
\end{minipage} & \begin{minipage}[b]{\linewidth}\raggedright
\(\approx\) time
\end{minipage} \\
\midrule\noalign{}
\endhead
\bottomrule\noalign{}
\endlastfoot
\textbf{Table 1} (unified benchmark) & \texttt{unified\_benchmark}, then
\texttt{scripts/plot\_unified\_panels.py} for the panel charts and
\texttt{\_tables.md} & 100 s \\
\textbf{Fig. 1} (order error vs the ACI bound) &
\texttt{order\_vs\_theory} & 30 s \\
§4 collapse statistics (\(\rho_S\), \(r\)) & derived from
\texttt{order\_vs\_theory.json}: rank and linear correlation of
\texttt{kendall} against \texttt{loss}, pooled over all four models and
all eleven levels & --- \\
\textbf{Figs. 2--3} (envelope; prefix pathology) & \texttt{choo\_bem} &
20 min \\
§5.3 recalibration & \texttt{recalibration} & 1.5 min \\
\textbf{Figs. 4, 5, 9} (payoff rule; \(r\)-sweep) &
\texttt{directional\_test} & 12 min \\
§5.4 decision-statistic variance & \texttt{measure\_payoff\_variance.py}
(not a \texttt{run\_} script) & 30 min \\
§5.5 eager-combiner mechanism check &
\texttt{tests/test\_combiner\_small.py} (console) & 2 s \\
\textbf{Fig. 6} (real predictor) & \texttt{real\_predictor} & 15 s \\
\textbf{Fig. 7} (six real graphs) \(\dagger\) &
\texttt{realworld\_robustness}, plus its \texttt{\_tables.md} & 65 s \\
§6.3 rank vs regression training & \texttt{rank\_vs\_mse\_mve},
\texttt{rank\_when\_it\_matters}, \texttt{rank\_real\_trace} & 40 s \\
\textbf{Fig. 8} (budget--stakes scissors) &
\texttt{impossibility\_frontier} & 6 s \\
§7 numerical verification & \texttt{scripts/verify\_witness\_gap.py} and
\texttt{verify\_budget\_stakes\_hetero.py} (console) & seconds \\
§8 serving case study & \texttt{serving}, \texttt{serving\_trace},
\texttt{serving\_dynamic}, \texttt{prefix\_cache},
\texttt{serving\_slo\_probe} & varies \\
Phase-2 reproduction of \cite{borodin2018experimental} &
\texttt{er\_full}, \texttt{left\_regular}, \texttt{realworld}
\(\dagger\) & 35 min \\
\end{longtable}
}

Scripts live under \texttt{scripts/} and write to \texttt{results/}; all
are run from the project root. Seven hand-verifiable correctness anchors
live under \texttt{tests/}
(\texttt{for\ t\ in\ tests/test\_*.py;\ do\ python3\ "\$t";\ done}).

\subsection{Order of execution}

Two dependencies are not self-evident: \texttt{plot\_unified\_panels.py}
and \texttt{run\_consistency\_robustness.py} both consume
\texttt{results/unified\_benchmark.json} and must run after
\texttt{run\_unified\_benchmark.py}. Everything else is independent.

\begin{verbatim}
# self-contained (clean checkout, no external data)
python3 scripts/run_unified_benchmark.py       # Table 1 (data)
python3 scripts/plot_unified_panels.py         # Table 1 (panel charts; needs the line above)
python3 scripts/run_order_vs_theory.py         # Fig 1
python3 scripts/run_directional_test.py        # Figs 4, 5, 9
python3 scripts/run_impossibility_frontier.py  # Fig 8
python3 scripts/verify_witness_gap.py          # Lemma 2, directional test, plug-in blindness
python3 scripts/verify_budget_stakes_hetero.py # heterogeneous profiles, slack identity
python3 scripts/measure_payoff_variance.py     # decision-statistic sd quoted in 5.4
python3 scripts/run_choo_bem.py                # Figs 2, 3   (~20 min)

# run on the traces shipped in data/trace/
python3 scripts/run_real_predictor.py          # Fig 6
python3 scripts/run_rank_real_trace.py         # 6.3, real-feature arm

# require data/realworld/ to be downloaded first (see A.4)
python3 scripts/run_realworld_robustness.py    # Fig 7
\end{verbatim}

\subsection{Data provenance}

The request traces ship with the artifact under \texttt{data/trace/}.
The third-party graph datasets do not, and must be placed under
\texttt{data/realworld/} before the daggered scripts will run.

\begin{itemize}
\tightlist
\item
  \textbf{Real graphs (Section 6.2).} Six graphs from the Network
  Repository (networkrepository.com): socfb-Caltech36, socfb-Reed98,
  bio-CE-GN, bio-CE-PG, econ-beause and econ-mbeaflw, downloaded as
  MatrixMarket or edge-list files, reduced to simple undirected graphs,
  and made bipartite by a random balanced partition.
\item
  \textbf{Traces (Sections 6.1, 6.3, 8).} Wikipedia ``top articles per
  day'' for four days from the Wikimedia REST pageviews API
  (\texttt{data/trace/wiki/}; the live day is the truth, earlier days
  are the 1-, 7- and 30-day-stale forecasts); the Azure LLM inference
  trace from Microsoft's public Azure dataset release
  (\texttt{data/trace/azure\_llm/}; context and generated token counts
  with timestamps); and the Mooncake conversation trace
  (\texttt{data/trace/mooncake/}; per-request prefix-cache block
  hashes). The Wikipedia data is a snapshot of a live source rather than
  a static benchmark, so exact reproduction requires the same snapshot,
  not merely the same API.
\end{itemize}

\section{Proof of Theorem 1(i)}

Section 7.5 gives the argument in outline; we record it in full here,
since it is the lower half of the paper's main theorem. Recall the
setting: a cell family with profile
\(\{(\theta_i,\varepsilon_i)\}_{i\le m}\), specialist mass
\(\sigma^2 = \sum_i\theta_i/N\) with
\(N=\sum_i(1+\theta_i)=\mathrm{OPT}\), and a scenario pair
\((G,\mathrm{Bd})\) that shares one advice and flips the
advice-agreement of a cell set \(W\) whose per-cell signals satisfy
\(\varepsilon_i \le \varepsilon_W\). Its payoff gap is
\(g=\tfrac2N\sum_{i\in W}\theta_i\varepsilon_i\).

\textbf{Claim.} Every test-and-fallback algorithm \(A_k\) that decides
by an arbitrary measurable (possibly randomized) function of the prefix
\(X_{1:k}\) and the advice, with \(k=o\!\big(1/(\varepsilon_W g)\big)\),
has \(\eta_c+\eta_r \ge 1-o(1)\).

\emph{Proof.} By Lemma 1 it suffices to show
\(\gamma_k := \mathrm{TV}(\mathcal L_G,\mathcal
L_{\mathrm{Bd}}) = o(1)\), where \(\mathcal L_\cdot\) is the law of the
prefix under each scenario; Lemma 1 then gives
\((1-\eta_c)\le\eta_r+\gamma_k+o(1)\), i.e. \(\eta_c+\eta_r\ge 1-o(1)\).

\emph{Step 1 (the coupling).} Arrivals are i.i.d., so
\(\mathcal L_G = P^{\otimes k}\) and
\(\mathcal L_{\mathrm{Bd}} = Q^{\otimes k}\) for the two per-arrival
laws \(P,Q\) on the type set. The two scenarios differ only in the
specialist bias \(s_i\) of the cells in \(W\): outside \(W\), and on
every flexible type, \(P\) and \(Q\) assign identical mass, and the
coupling that leaves those arrivals untouched is exact there. Inside a
flipped cell \(i\in W\), the two specialist types carry masses
\(\tfrac{\theta_i}{N}(\tfrac12+\varepsilon_i)\) and
\(\tfrac{\theta_i}{N}(\tfrac12-\varepsilon_i)\) under \(P\), and the
same two masses swapped under \(Q\).

\emph{Step 2 (per-sample Hellinger distance).} Writing
\(H^2(P,Q)=\sum_x(\sqrt{P(x)}-\sqrt{Q(x)})^2\), only the \(2|W|\)
swapped atoms contribute, and cell \(i\) contributes
\[2\,\frac{\theta_i}{N}\Big(\sqrt{\tfrac12+\varepsilon_i}-\sqrt{\tfrac12-\varepsilon_i}\Big)^{2}
\;=\; \frac{2\theta_i}{N}\Big(1-\sqrt{1-4\varepsilon_i^{2}}\Big),\]
using \((\sqrt a-\sqrt b)^2 = a+b-2\sqrt{ab}\) with \(a+b=1\) and
\(ab=\tfrac14-\varepsilon_i^2\). Since \(2t \le 1-\sqrt{1-4t} \le 4t\)
for \(t=\varepsilon_i^2\in[0,\tfrac14]\),
\[H^2(P,Q) \;=\; \sum_{i\in W}\frac{2\theta_i}{N}\Big(1-\sqrt{1-4\varepsilon_i^{2}}\Big)
\;\in\; [4,8]\cdot\frac{1}{N}\sum_{i\in W}\theta_i\varepsilon_i^{2}.\]
Bounding each \(\varepsilon_i\) by \(\varepsilon_W\) in one factor,
\[H^2(P,Q) \;\le\; \frac{8}{N}\sum_{i\in W}\theta_i\varepsilon_i^{2}
\;\le\; 4\varepsilon_W\cdot\frac{2}{N}\sum_{i\in W}\theta_i\varepsilon_i
\;=\; 4\,\varepsilon_W\,g .\]

\emph{Step 3 (tensorization).} With
\(h^2 := H^2/2 = 1-\mathrm{BC}(P,Q)\) the normalized squared Hellinger
distance, the Bhattacharyya coefficient of a product measure is the
product of the coefficients, so
\(1-h^2(P^{\otimes k},Q^{\otimes k}) = (1-h^2(P,Q))^k\) and hence
\(h^2(P^{\otimes k},Q^{\otimes k}) \le k\,h^2(P,Q)\).

\emph{Step 4 (to total variation).} The standard comparison
\(\mathrm{TV}\le h\sqrt{2-h^{2}}\le\sqrt2\,h\) gives
\[\gamma_k \;\le\; \sqrt2\,h(P^{\otimes k},Q^{\otimes k}) \;\le\; \sqrt{2k\,h^2(P,Q)}
\;=\; \sqrt{k\,H^2(P,Q)} \;\le\; \sqrt{4k\,\varepsilon_W\,g}.\] Steps
3--4, chained with Lemma 1, are machine-checked in the form
\((1-\eta_c)\le\eta_r+\sqrt{2k\,(1-\mathrm{BC}(P,Q))}\) for arbitrary
finite \(P,Q\) and every \([0,1]\)-valued rule on the \(k\)-prefix
(\texttt{Scenario.master\_tradeoff\_iid}, Appendix C); Steps 1--2 are
the explicit evaluation on the family.

\emph{Step 5 (conclusion).} If \(k=o\!\big(1/(\varepsilon_W g)\big)\)
then \(k\,\varepsilon_W g\to0\), so \(\gamma_k\to0\): the prefix is
asymptotically uninformative about which scenario is in force, and Lemma
1 converts that into \(\eta_c+\eta_r\ge1-o(1)\). Here \(o(1)\) is taken
along any sequence of families and prefix lengths with
\(k\varepsilon_W g\to0\); the \(o(1)\) inside Lemma 1 additionally
absorbs the \(O(k/n)\) contribution of the prefix itself to the ratio,
which is why the sublinearity assumption \(k=o(n)\) is in force.
\(\qed\)

\textbf{Where the two halves meet.} Theorem 1(ii) spends
\(k=O(\sigma^2/\min(\delta,\Delta)^2)\), while (i) forbids
\(k=o(1/(\varepsilon_W g))\). On a balanced pair
\(\min(\delta,\Delta)=g/2\), so the gap between the two is the ratio
\(\sigma^2\varepsilon_W/g\). Cauchy--Schwarz on the definition of \(g\)
gives
\(g^2 \le \tfrac4N\,\sigma_W^2\sum_{i\in W}\theta_i\varepsilon_i^2\)
with \(\sigma_W^2\) the flipped mass, so that ratio is \(\Theta(1)\) ---
and the two sides agree up to logarithms --- exactly when the stakes are
carried, at comparable signal levels, by a constant fraction of the
specialist mass. When instead they hide in a vanishing sliver of
low-signal cells the ratio can diverge; that regime is open (Section
7.8).

\section{Machine-checked statements}

The statistical chain behind Section 7 is formalized in Lean 4 (v4.33.0)
over Mathlib, in the \texttt{budgetstakes/} directory of the artifact:
seven files, 66 theorems and lemmas, no \texttt{sorry}.
\texttt{lake\ build} replays the check in about a minute once the
Mathlib cache is fetched, and the repository's continuous integration
runs it on every commit. The development is deliberately self-contained
--- a distribution is a vector of real weights on a finite type,
expectations are finite sums, and independence enters through a single
combinatorial identity,
\(\sum_{x\in\Omega^k}\prod_i g(x_i)=\big(\sum_\omega g(\omega)\big)^k\)
--- so that a reader can audit it without measure theory. The table maps
the paper's statements to theorem names.

{ % do not increment counter
\begin{longtable}[]{@{}
  >{\raggedright\arraybackslash}p{(\linewidth - 4\tabcolsep) * \real{0.2308}}
  >{\raggedright\arraybackslash}p{(\linewidth - 4\tabcolsep) * \real{0.3846}}
  >{\raggedright\arraybackslash}p{(\linewidth - 4\tabcolsep) * \real{0.3846}}@{}}
\toprule\noalign{}
\begin{minipage}[b]{\linewidth}\raggedright
Statement in the paper
\end{minipage} & \begin{minipage}[b]{\linewidth}\raggedright
Lean theorem
\end{minipage} & \begin{minipage}[b]{\linewidth}\raggedright
Form proved
\end{minipage} \\
\midrule\noalign{}
\endhead
\bottomrule\noalign{}
\endlastfoot
Lemma 2 (payoff identity), §7.4 & \texttt{CellFamily.payoff\_identity} &
heterogeneous profile, arbitrary truth biases, from the cell
constants \\
\(1-\rho_{\mathrm{base}}=\sigma^2/2\), §7.3 &
\texttt{CellFamily.}\allowbreak\texttt{slack\_eq\_half\_specialistMass}
& exact \\
affine law and
\(\rho_{\mathrm{perfect}}-\rho_{\mathrm{base}}\le 2\varepsilon_{\max}(1-\rho_{\mathrm{base}})\),
§7.3 & \texttt{CellFamily.affine\_law},
\texttt{CellFamily.}\allowbreak\texttt{rhoPerfect\_sub\_rhoBase\_le\_slack}
& exact \\
Lemma 1 (master trade-off), §7.2 & \texttt{Scenario.master\_tradeoff} &
exact, no \(o(1)\): \((1-\eta_c)\le\eta_r+\mathrm{TV}\) for every
\([0,1]\)-valued rule on a finite prefix space \\
Theorem 1(ii), upper half &
\texttt{FinDist.}\allowbreak\texttt{iid\_prob\_sum\_nonpos\_le} & finite
Chernoff, \(P^{\otimes k}(\sum c\le0)\le e^{-k\mu^2/4}\) for
\(\lvert c\rvert\le1\), \(\mu=\mathbb E[c]\); with
\(\mu\ge2\min(\delta,\Delta)\) this is error
\(\le e^{-k\min(\delta,\Delta)^2}\) \\
Theorem 1(ii), \(\sigma^2\)-scaling &
\texttt{FinDist.}\allowbreak\texttt{iid\_prob\_sum\_nonpos\_le\_cheb} &
exact prefix variance \(k\,\mathrm{Var}(c)\) and Chebyshev, error
\(\le\mathrm{Var}(c)/(k\mu^2)\); with
\(\mathrm{Var}(c)\le\mathbb E[c^2]=\sigma^2\) this is
\(\sigma^2/(4k\min(\delta,\Delta)^2)\) \\
Theorem 1(i): Appendix B Steps 3--4 chained with Lemma 1 &
\texttt{Scenario.}\allowbreak\texttt{master\_tradeoff\_iid} &
\((1-\eta_c)\le\eta_r+\sqrt{2k\,(1-\mathrm{BC}(P,Q))}\) for arbitrary
finite \(P,Q\) and every rule on the \(k\)-prefix (tensorization
\texttt{bc\_iid}, comparison \texttt{tvDist\_iid\_le}) \\
Remark in §7.8 (random order), lower half &
\texttt{FinDist.}\allowbreak\texttt{shuffled\_iid\_expect} & a uniformly
shuffled i.i.d. sample has the product law \\
Remark in §7.8 (random order), upper half &
\texttt{RandomOrder.}\allowbreak\texttt{ro\_prob\_sum\_nonpos\_le} &
without-replacement prefix sum: variance
\(\le k\,\mathbb E_{\mathrm{pop}}[c^2]\), error
\(\le\mathbb E_{\mathrm{pop}}[c^2]/(k\mu^2)\) \\
\end{longtable}
}

What is \emph{not} formalized, so that the certificate is not over-read:
(a) the cell constants of §7.3 (\(\mathrm{OPT}=1+\theta_i\), baseline
\(1+\theta_i/2\), follow gain \(\theta_i\lvert s_i-\tfrac12\rvert\)),
which the development takes as the definition of the model; (b) Steps
1--2 of Appendix B --- the explicit evaluation
\(1-\mathrm{BC}=\tfrac1N\sum_{i\in W}\theta_i\big(1-\sqrt{1-4\varepsilon_i^2}\big)\)
and the \([2t,4t]\) bounds --- two lines of algebra on the page; (c)
Bernstein's inequality, i.e.~the \(\sigma^2\)-dependence \emph{inside
the exponent} of Theorem 1(ii) (the \(\sigma^2\) scaling is certified at
Chebyshev grade only); (d) the \(o(1)\) bookkeeping of the prefix's own
\(O(k/n)\) contribution, which the formal statements avoid by defining
values as post-decision ratios; and (e) exponential tails under random
order, which we cite rather than reprove. Nothing about the matching
algorithms themselves is formalized: the development is about the
statistics of testing advice, which is what the law is.
